\chapter{Aritmetika: Pembagi dan Bilangan Prima}\label{ch:g9:arith}

Aritmetika mempelajari bilangan bulat dan cara bilangan yang satu
membagi yang lain. Tokoh utamanya adalah \omterm{def:g9:arith:prime}{bilangan prima}, yaitu balok
pembangun yang darinya setiap bilangan bulat disusun lewat perkalian.
Bab ini berakhir dengan \omterm{def:g9:arith:gcd}{faktor persekutuan terbesar}, alat yang tepat
untuk menyederhanakan \omterm{def:g9:fractions:fraction}{pecahan} sekali untuk selamanya. Kisah ini
berlanjut, jauh lebih jauh, di buku Sekolah Menengah dan sesudahnya.

\section{Pembagi dan kelipatan}

\begin{definition}[Pembagi, kelipatan]\label{def:g9:arith:divisor}
Misalkan $a$ dan $b$ bilangan bulat positif. Kita katakan bahwa $b$
\emph{membagi}\index{pembagi} $a$ (atau bahwa $b$ pembagi $a$, atau
bahwa $a$ adalah \emph{kelipatan}\index{kelipatan} $b$) ketika
$a = b \times k$ untuk suatu bilangan bulat $k$ --- yaitu ketika
pembagian $a$ oleh $b$ meninggalkan \omterm{def:g3:division:remainder}{sisa} $0$.
\end{definition}

\begin{example}
\omterm{def:g9:arith:divisor}{Pembagi} $24$ adalah $1, 2, 3, 4, 6, 8, 12, 24$ --- semuanya datang
berpasangan yang \omterm{def:g2:mult:def}{hasil kalinya} $24$: $(1,24)$, $(2,12)$, $(3,8)$,
$(4,6)$. \omterm{def:g5:division:multiple}{Kelipatan} $7$ adalah $7, 14, 21, 28, \dots$
\end{example}

\begin{proposition}[Aturan keterbagian]\label{prop:g9:arith:rules}
Sebuah bilangan bulat \omterm{def:g6:wholes:divisible}{habis dibagi}:
\begin{itemize}
  \item $2$ ketika angka terakhirnya \omterm{def:g3:numbers:evenodd}{genap} ($0, 2, 4, 6, 8$);
  \item $5$ ketika angka terakhirnya $0$ atau $5$;
  \item $10$ ketika angka terakhirnya $0$;
  \item $3$ (berturut-turut $9$) ketika \omterm{def:g1:addition:def}{jumlah} angkanya \omterm{def:g6:wholes:divisible}{habis dibagi}
        $3$ (berturut-turut $9$);
  \item $4$ ketika dua angka terakhirnya membentuk bilangan yang \omterm{def:g6:wholes:divisible}{habis dibagi} $4$.
\end{itemize}
\end{proposition}
\admitted

\begin{example}
$7\,215$ berakhir dengan $5$: jadi \omterm{def:g6:wholes:divisible}{habis dibagi} $5$. \omterm{def:g1:addition:def}{Jumlah} angkanya
$7 + 2 + 1 + 5 = 15$, yang \omterm{def:g6:wholes:divisible}{habis dibagi} $3$ tetapi tidak \omterm{def:g6:wholes:divisible}{habis dibagi}
$9$: jadi $7\,215$ \omterm{def:g6:wholes:divisible}{habis dibagi} $3$, tidak \omterm{def:g6:wholes:divisible}{habis dibagi} $9$. Memang
$7\,215 = 3 \times 5 \times 481$.
\end{example}

\section{Bilangan prima}

\begin{definition}[Bilangan prima]\label{def:g9:arith:prime}
\emph{Bilangan prima}\index{bilangan prima} adalah bilangan bulat
$\geq 2$ yang \omterm{def:g9:arith:divisor}{pembaginya} hanya $1$ dan dirinya sendiri. Bilangan prima
di bawah $30$ adalah
\[
2,\ 3,\ 5,\ 7,\ 11,\ 13,\ 17,\ 19,\ 23,\ 29 .
\]
Bilangan $1$ \emph{bukan} prima (menurut kesepakatan), dan bilangan
bulat $\geq 2$ yang bukan prima disebut \emph{komposit}.
\end{definition}

\begin{theorem}[Pemfaktoran prima]\label{thm:g9:arith:factorization}
Setiap bilangan bulat $\geq 2$ adalah \omterm{def:g2:mult:def}{hasil kali} \omterm{def:g9:arith:prime}{bilangan prima}, dan
pemfaktoran itu tunggal kecuali urutan faktornya.
\end{theorem}
\admitted

\begin{method}[Memfaktorkan bilangan bulat]\label{met:g9:arith:factorization}
Bagilah dengan \omterm{def:g9:arith:prime}{bilangan prima} terkecil yang mungkin, berulang-ulang,
sampai mencapai $1$:
\begin{enumerate}
  \item cobalah $2$ selama bilangannya \omterm{def:g3:numbers:evenodd}{genap};
  \item lalu cobalah $3$, lalu $5$, lalu $7$, \dots\ (hanya yang prima);
  \item berhentilah ketika \omterm{def:g3:division:remainder}{hasil baginya} $1$; kumpulkan faktornya
        beserta \omterm{def:g8:powers:def}{eksponennya}.
\end{enumerate}
Cukup mencoba \omterm{def:g9:arith:prime}{bilangan prima} $p$ dengan $p^2$ tidak melampaui bilangan
yang sedang dikerjakan: kalau tidak ada yang membaginya, maka
bilangannya sendiri prima.
\end{method}

\begin{example}\label{ex:g9:arith:factorization}
Faktorkan $360$, satu pembagian setiap kalinya:
\[
360 = 2 \times 180, \quad
180 = 2 \times 90, \quad
90 = 2 \times 45, \quad
45 = 3 \times 15, \quad
15 = 3 \times 5,
\]
jadi
\[
360 = 2 \times 2 \times 2 \times 3 \times 3 \times 5 = 2^3 \times 3^2
\times 5 .
\]
\end{example}

\begin{omfigure}
\begin{tikzpicture}[scale=0.78]
  \node (n360) at (0,0) {$360$};
  \node (n180) at (1.2,-1) {$180$};
  \node (n90) at (2.4,-2) {$90$};
  \node (n45) at (3.6,-3) {$45$};
  \node (n15) at (4.8,-4) {$15$};
  \node (n5) at (6.0,-5) {$5$};
  \node[omThm] (p2a) at (-1.2,-1) {$2$};
  \node[omThm] (p2b) at (0,-2) {$2$};
  \node[omThm] (p2c) at (1.2,-3) {$2$};
  \node[omThm] (p3a) at (2.4,-4) {$3$};
  \node[omThm] (p3b) at (3.6,-5) {$3$};
  \draw (n360) -- (n180); \draw (n360) -- (p2a);
  \draw (n180) -- (n90); \draw (n180) -- (p2b);
  \draw (n90) -- (n45); \draw (n90) -- (p2c);
  \draw (n45) -- (n15); \draw (n45) -- (p3a);
  \draw (n15) -- (n5); \draw (n15) -- (p3b);
\end{tikzpicture}

{\small Pohon faktor $360$: tiap langkahnya memisahkan faktor prima
terkecilnya (berwarna merah). Membaca daun merahnya dan $5$ yang
terakhir: $360 = 2^3 \times 3^2 \times 5$.}
\end{omfigure}

\begin{theorem}[Euclid]\label{thm:g9:arith:euclid}
\omterm{def:g9:arith:prime}{Bilangan prima} ada tak berhingga banyaknya.
\end{theorem}

\begin{proof}
Andaikan banyaknya berhingga saja, katakanlah $p_1, p_2, \dots, p_k$,
lalu tinjaulah
\[
N = p_1 \times p_2 \times \dots \times p_k + 1 .
\]
Membagi $N$ dengan $p_i$ mana pun meninggalkan \omterm{def:g3:division:remainder}{sisa} $1$, jadi tidak
ada $p_i$ yang membagi $N$. Tetapi $N \geq 2$ punya sedikitnya satu
\omterm{def:g9:arith:divisor}{pembagi} prima (\cref{thm:g9:arith:factorization}) --- yaitu \omterm{def:g9:arith:prime}{bilangan
prima} yang tidak ada dalam daftar kita. Kontradiksi: tidak ada daftar
berhingga yang dapat memuat semua \omterm{def:g9:arith:prime}{bilangan primanya}.
\end{proof}

\section{Faktor persekutuan terbesar}

\begin{definition}[FPB]\label{def:g9:arith:gcd}
\emph{Faktor persekutuan terbesar}\index{faktor persekutuan terbesar}
dua bilangan bulat positif $a$ dan $b$, yang ditulis $\gcd(a, b)$,
adalah bilangan bulat terbesar yang membagi keduanya. Ketika
$\gcd(a, b) = 1$, kedua bilangan bulatnya disebut \emph{saling
prima}\index{saling prima}: keduanya tidak berbagi \omterm{def:g9:arith:divisor}{pembagi} selain $1$.
\end{definition}

\begin{example}
\omterm{def:g9:arith:divisor}{Pembagi} $18$: $1, 2, 3, 6, 9, 18$. \omterm{def:g9:arith:divisor}{Pembagi} $24$: $1, 2, 3, 4, 6,
8, 12, 24$. \omterm{def:g9:arith:divisor}{Pembagi} bersamanya: $1, 2, 3, 6$; jadi
$\gcd(18, 24) = 6$. Bilangan bulat $15$ dan $28$ \omterm{def:g9:arith:gcd}{saling prima}.
\end{example}

\begin{proposition}[FPB dari pemfaktorannya]\label{prop:g9:arith:gcdfact}
FPB dua bilangan bulat adalah \omterm{def:g2:mult:def}{hasil kali} \omterm{def:g9:arith:prime}{bilangan prima} yang muncul
pada \emph{kedua} pemfaktorannya, masing-masing diambil dengan
\omterm{def:g8:powers:def}{eksponen} yang \emph{lebih kecil} dari kedua \omterm{def:g8:powers:def}{eksponennya}.
\end{proposition}
\admitted

\begin{example}\label{ex:g9:arith:gcdfact}
$360 = 2^3 \times 3^2 \times 5$ dan $84 = 2^2 \times 3 \times 7$.
\omterm{def:g9:arith:prime}{Bilangan prima} bersamanya: $2$ (\omterm{def:g8:powers:def}{eksponennya} $3$ dan $2$: ambil $2$)
dan $3$ (\omterm{def:g8:powers:def}{eksponennya} $2$ dan $1$: ambil $1$). Jadi
\[
\gcd(360, 84) = 2^2 \times 3 = 12 .
\]
\end{example}

\begin{theorem}[Algoritma Euclid]\label{thm:g9:arith:euclidalgo}
Kalau $a = bq + r$ adalah pembagian $a$ oleh $b$ dengan \omterm{def:g3:division:remainder}{sisa} $r$, maka
\[
\gcd(a, b) = \gcd(b, r).
\]
Dengan mengulang pembagiannya sampai \omterm{def:g3:division:remainder}{sisanya} $0$, FPB dari $a$ dan $b$
adalah \emph{\omterm{def:g3:division:remainder}{sisa} terakhir yang bukan \omterm{def:g1:counting:zero}{nol}}.
\end{theorem}

\begin{proof}
Dari $a = bq + r$: setiap bilangan bulat yang membagi $b$ dan $r$
membagi $bq + r = a$; dan dari $r = a - bq$: setiap bilangan bulat
yang membagi $a$ dan $b$ membagi $r$. Jadi pasangan $(a, b)$ dan
$(b, r)$ punya \omterm{def:g9:arith:divisor}{pembagi} bersama yang persis sama --- khususnya yang
terbesarnya sama. Karena \omterm{def:g3:division:remainder}{sisanya} menurun secara tegas, algoritmanya
berhenti, dan $\gcd(x, 0) = x$ memberi \omterm{def:g3:division:remainder}{sisa} terakhir yang bukan \omterm{def:g1:counting:zero}{nol}.
\end{proof}

\begin{example}\label{ex:g9:arith:euclidalgo}
Hitunglah $\gcd(1071, 462)$:
\begin{align*}
1071 &= 462 \times 2 + 147, \\
462 &= 147 \times 3 + 21, \\
147 &= 21 \times 7 + 0 .
\end{align*}
\omterm{def:g3:division:remainder}{Sisa} terakhir yang bukan \omterm{def:g1:counting:zero}{nol} adalah $21$: $\gcd(1071, 462) = 21$.
\end{example}

\begin{method}[Menyederhanakan pecahan sepenuhnya]\label{met:g9:arith:simplify}
Untuk menulis $\dfrac ab$ sesederhana mungkin:
\begin{enumerate}
  \item hitunglah $d = \gcd(a, b)$, misalnya dengan algoritma Euclid;
  \item bagilah \omterm{def:g4:fractions:def}{pembilang} dan \omterm{def:g4:fractions:def}{penyebutnya} dengan $d$:
        $\dfrac ab = \dfrac{a \div d}{b \div d}$;
  \item \omterm{def:g9:fractions:fraction}{pecahan} yang dihasilkan bersifat \emph{tak tersederhanakan}:
        \omterm{def:g4:fractions:def}{pembilang} dan \omterm{def:g4:fractions:def}{penyebutnya} \omterm{def:g9:arith:gcd}{saling prima}.
\end{enumerate}
\end{method}

\begin{example}
$\dfrac{462}{1071} = \dfrac{462 \div 21}{1071 \div 21} = \dfrac{22}{51}$,
dan $\gcd(22, 51) = 1$: jadi tak tersederhanakan.
\end{example}

\section{Latihan}

\begin{exercise}[$\star$]\label{exo:g9:arith:1}
Sebutkan semua \omterm{def:g9:arith:divisor}{pembagi} $36$, \omterm{def:g9:arith:divisor}{pembagi} $45$, dan \omterm{def:g9:arith:divisor}{pembagi} $17$.
\end{exercise}

\begin{exercise}[$\star$]\label{exo:g9:arith:2}
Dengan memakai aturan keterbagiannya, tentukan apakah $2\,346$ \omterm{def:g6:wholes:divisible}{habis
dibagi} $2$, $3$, $4$, $5$, dan $9$.
\end{exercise}

\begin{exercise}[$\star$]\label{exo:g9:arith:3}
Berikan pemfaktoran prima $72$, $150$, $210$ dan $121$.
\end{exercise}

\begin{exercise}[$\star$]\label{exo:g9:arith:4}
Apakah $101$ prima? Apakah $91$? Apakah $143$? Berilah alasan dengan
memakai aturan berhenti pada \cref{met:g9:arith:factorization}.
\end{exercise}

\begin{exercise}[$\star$]\label{exo:g9:arith:5}
Hitunglah $\gcd(48, 60)$ dengan dua cara: dengan menyebutkan \omterm{def:g9:arith:divisor}{pembagi}
bersamanya, dan dari pemfaktoran primanya.
\end{exercise}

\begin{exercise}[$\star\star$]\label{exo:g9:arith:6}
Pakailah algoritma Euclid untuk menghitung $\gcd(255, 154)$, lalu
$\gcd(1053, 325)$. Tulislah tiap baris pembagiannya.
\end{exercise}

\begin{exercise}[$\star\star$]\label{exo:g9:arith:7}
Jadikan \omterm{def:g9:fractions:fraction}{pecahan} $\dfrac{588}{504}$ tak tersederhanakan. (Hitunglah
FPB-nya dengan cara pilihanmu, lalu bagilah.)
\end{exercise}

\begin{exercise}[$\star\star$]\label{exo:g9:arith:8}
Seorang penjual bunga punya $84$ mawar dan $126$ tulip lalu ingin
membuat rangkaian yang identik, dengan memakai semua bunganya, dan
dengan rangkaian sebanyak mungkin. Berapa rangkaian yang dapat
dibuatnya, dan apa isi masing-masingnya?
\end{exercise}

\begin{exercise}[$\star\star$]\label{exo:g9:arith:9}
Dua kapal feri berangkat dari dermaga yang sama pukul 8:00. Yang satu
berangkat tiap $24$ menit, yang lain tiap $36$ menit. Pukul berapa
keduanya berikutnya berangkat bersama-sama? (Carilah \omterm{def:g9:arith:divisor}{kelipatan}
bersama terkecil dari $24$ dan $36$; pemfaktorannya membantu.)
\end{exercise}

\begin{exercise}[$\star\star\star$]\label{exo:g9:arith:10}
Misalkan $n$ sebuah bilangan bulat positif.
\begin{enumerate}
  \item Tunjukkan bahwa $\gcd(n, n+1) = 1$ (bilangan bulat yang
        berurutan selalu \omterm{def:g9:arith:gcd}{saling prima}).
  \item Simpulkan bahwa \omterm{def:g9:fractions:fraction}{pecahan} $\dfrac{n}{n+1}$ selalu tak tersederhanakan.
\end{enumerate}
\end{exercise}

\section{Soal: Kendi air, jangkrik dan seratus loker}

\begin{problem}[{Soal akhir pekan --- FPB memutuskan takaran mana
yang dapat diukur dua kendi; bilangan prima melindungi jangkrik; dan
loker yang tetap terbuka adalah kuadrat sempurna}]
\label{pb:g9:arith:1}
Tiga teka-teki yang tampak seperti tebakan padahal sungguh-sungguh
aritmetika: menakar air dengan kendi tanpa tanda ukur (yaitu FPB dalam
samaran), daur hidup serangga yang berkembang menjadi \omterm{def:g9:arith:prime}{bilangan prima},
dan sebuah lorong terkenal berisi seratus loker yang keadaan akhirnya
ditentukan dengan membilang \omterm{def:g9:arith:divisor}{pembagi}. Semuanya berjalan dengan mesin
bab ini: keterbagian, pemfaktoran prima
(\cref{thm:g9:arith:factorization}) dan algoritma Euclid
(\cref{thm:g9:arith:euclidalgo}).

\medskip
\noindent\textbf{Bagian I --- Kendi airnya.}
Kamu berdiri di sebuah pancuran dengan dua kendi tanpa tanda ukur,
berukuran $5$ L dan $3$ L. Gerakan yang diizinkan: mengisi sebuah
kendi sampai penuh, mengosongkan sebuah kendi sepenuhnya, menuang satu
kendi ke kendi yang lain sampai sumbernya kosong atau tujuannya penuh.
\begin{enumerate}
  \item Takarlah tepat $1$ L. (Uraikan urutan gerakanmu
        dan isi kedua kendinya sesudah masing-masingnya.)
  \item Takarlah tepat $4$ L --- yaitu teka-teki dari sebuah film
        laga yang terkenal. (Hal itu dapat dilakukan dalam enam gerakan.)
  \item Bilangan bulat liter yang mana dari $1$ sampai $8$ yang dapat
        kamu tunjukkan (dalam satu kendi, atau terbagi pada
        keduanya)? Lengkapilah daftarnya, dengan memakai ulang
        urutanmu.
  \item Kendi yang baru: $6$ L dan $4$ L. Cobalah menakar $1$ L ---
        lalu jelaskan mengapa hal itu sia-sia: periksalah bahwa
        masing-masing dari ketiga gerakan yang diizinkan menjaga isi
        tiap kendinya tetap \omterm{def:g9:arith:divisor}{kelipatan}
        $2$, jadi setiap takaran yang tercapai bernilai \omterm{def:g3:numbers:evenodd}{genap}.
  \item Alasan pada pertanyaan 4 berlaku umum: dengan kendi $a$
        dan $b$ liter, setiap takaran yang tercapai adalah \omterm{def:g9:arith:divisor}{kelipatan}
        $\gcd(a, b)$. Hitunglah $\gcd(6, 4)$ dan $\gcd(5, 3)$, lalu
        katakan apa yang diramalkan hukumnya untuk tiap pasangan kendinya.
\end{enumerate}

\medskip
\noindent\textbf{Bagian II --- Euclid di pancuran.}
\begin{enumerate}[resume]
  \item Hitunglah dengan algoritma Euclid: $\gcd(91, 65)$ dan
        $\gcd(2\,026, 46)$.
  \item Jelaskan, dengan kata-katamu sendiri, mengapa takaran yang
        muncul pada kendinya adalah \omterm{def:g3:division:remainder}{sisa} milik Euclid dalam samaran:
        dengan kendi $13$ L dan $5$ L, isilah kendi kecilnya
        berulang-ulang lalu tuangkan ke kendi besarnya (dengan
        mengosongkan kendi besarnya setiap kali ia penuh). Takaran
        baru yang mana yang muncul lebih dahulu --- lalu bandingkan
        semuanya dengan \omterm{def:g3:division:remainder}{sisa} pada algoritma Euclid
        untuk $(13, 5)$.
  \item Simpulkan jawaban sang juara: dengan kendi $13$ dan $5$
        liter, dapatkah kamu menakar tepat $1$ L? Berilah alasan
        dalam satu baris dengan pertanyaan 5 dan $\gcd(13, 5)$.
  \item Bukti kilat tentang \omterm{def:g9:arith:gcd}{saling prima} dengan gaya
        \cref{exo:g9:arith:10}: tunjukkan bahwa $\gcd(n, 2n + 1) = 1$
        untuk setiap bilangan bulat positif $n$. (Apa yang harus
        dibagi oleh \omterm{def:g9:arith:divisor}{pembagi} bersama $n$ dan $2n + 1$?)
  \item Dua bus berangkat bersama dari terminal pukul 7:00; yang satu
        berangkat tiap $12$ menit, yang lain tiap $18$. Sebutkan
        waktu keberangkatan berikutnya untuk masing-masingnya lalu
        carilah saat pertama keduanya berangkat bersama lagi.
        Periksalah pada contoh ini hukum yang indah: (\omterm{def:g9:arith:divisor}{kelipatan}
        bersama pertamanya) $\times$
        $\gcd$ $=$ \omterm{def:g2:mult:def}{hasil kali} kedua bilangannya --- lalu ujilah
        sekali lagi pada $5$ dan $3$.
\end{enumerate}

\medskip
\noindent\textbf{Bagian III --- Jangkrik, \omterm{def:g9:arith:divisor}{pembagi} dan loker.}
\begin{enumerate}[resume]
  \item Jangkrik Amerika Utara tertentu keluar hanya tiap $17$
        tahun; andaikan populasi pemangsanya memuncak tiap $4$
        tahun. Kalau keduanya terjadi tahun ini, berapa tahun lagi
        kemunculan jangkriknya berikutnya berbarengan dengan
        puncaknya? Pertanyaan yang sama kalau daur jangkriknya
        $16$ tahun --- seberapa sering mereka lalu dibantai?
        Jelaskan dalam satu kalimat mengapa evolusi mendorong
        daurnya ke panjang yang \emph{prima}.
  \item Dengan memakai pemfaktoran $360 = 2^3 \times 3^2 \times 5$,
        bilanglah \omterm{def:g9:arith:divisor}{pembagi} $360$ tanpa menyebutkannya satu per satu:
        sebuah \omterm{def:g9:arith:divisor}{pembagi} memilih \omterm{def:g8:powers:def}{eksponen} untuk $2$ (empat pilihan:
        $0, 1, 2, 3$), satu untuk $3$, satu untuk $5$. Berapa
        \omterm{def:g9:arith:divisor}{pembagi} seluruhnya?
  \item Tunjukkan bahwa pada pemfaktoran kuadrat sempurna
        $n = m^2$, setiap \omterm{def:g9:arith:prime}{bilangan primanya} membawa \omterm{def:g8:powers:def}{eksponen} yang
        \emph{\omterm{def:g3:numbers:evenodd}{genap}}. Simpulkan, tanpa menghitung akar kuadrat apa
        pun, bahwa $360$ bukan kuadrat sempurna.
  \item Pasangkan tiap \omterm{def:g9:arith:divisor}{pembagi} $d$ dari $n$ dengan pasangannya
        $\frac{n}{d}$ (untuk $n = 36$: $1 \leftrightarrow 36$,
        $2 \leftrightarrow 18$, $3 \leftrightarrow 12$,
        $4 \leftrightarrow 9$, $6 \leftrightarrow 6$). Kapan sebuah
        \omterm{def:g9:arith:divisor}{pembagi} menjadi pasangannya sendiri? Simpulkan kriterianya:
        $n$ punya banyak \omterm{def:g9:arith:divisor}{pembagi} yang \emph{\omterm{def:g3:numbers:evenodd}{ganjil}} persis ketika
        $n$ adalah kuadrat sempurna. Periksalah pada $36$ dan pada $360$.
  \item Seratus lokernya. Loker $1$ sampai $100$ bermula tertutup.
        Murid $1$ membalik keadaan setiap loker; murid $2$ membalik
        loker $2, 4, 6, \dots$; murid $k$ membalik
        \omterm{def:g9:arith:divisor}{kelipatan} $k$; dan seterusnya sampai murid $100$. Jelaskan
        murid mana saja yang menyentuh loker $n$, berapa kali loker
        itu terbalik keadaannya, dan --- dengan memakai pertanyaan
        14 --- loker mana persisnya yang berakhir terbuka. Berapa
        banyak yang terbuka?

\end{enumerate}
\end{problem}
